\documentclass[letterpaper,10pt]{article}
% =============================================================================
% Pathways: Economics & Finance Workshop — Document Preamble
% =============================================================================
% Usage: Each handout/activity .tex file should start with:
%   \documentclass[letterpaper,10pt]{article}
%   % =============================================================================
% Pathways: Economics & Finance Workshop — Document Preamble
% =============================================================================
% Usage: Each handout/activity .tex file should start with:
%   \documentclass[letterpaper,10pt]{article}
%   % =============================================================================
% Pathways: Economics & Finance Workshop — Document Preamble
% =============================================================================
% Usage: Each handout/activity .tex file should start with:
%   \documentclass[letterpaper,10pt]{article}
%   \input{../../codes/Preambles/header_doc}
% Compile: cd output/Handouts && TEXINPUTS=../../codes/Preambles:$TEXINPUTS xelatex file.tex
% =============================================================================

% ---------- Packages ----------
% Fonts
\usepackage{palatino}
\usepackage[utf8]{inputenc}
\usepackage[T1]{fontenc}
\usepackage[normalem]{ulem}

% Math
\usepackage{amsmath,amsthm}
\usepackage{newpxmath}
\usepackage{bm}
\usepackage{mathrsfs}
\usepackage{dsfont}

% Figures and tables
\usepackage{graphicx}
\usepackage{placeins}
\usepackage{xcolor}
\usepackage{array,tabularx,booktabs,longtable,subcaption}
\usepackage{threeparttable}
\usepackage{pdflscape}
\usepackage{siunitx}

% TikZ and pgfplots
\usepackage{pgfplots}
\pgfplotsset{compat=1.18}
\usepackage{tikz}
\usetikzlibrary{positioning,arrows.meta,calc}

% Colored boxes
\usepackage[most]{tcolorbox}

% Citations (natbib loaded before hyperref to avoid conflicts)
\usepackage{natbib}

% Miscellaneous
\usepackage{hyperref}
% Links stay clickable but render as plain text: `hidelinks` removes BOTH the
% coloured border boxes (hyperref's default for links/citations/URLs) and any
% link colouring. Must come after the package load.
\hypersetup{hidelinks}
\usepackage{setspace}
\usepackage{enumitem}

% ---------- Margins ----------
\usepackage[left=1in,right=1in,top=0.75in,bottom=1in,marginparwidth=1.5in,footskip=0.7in]{geometry}

% ---------- Itemize ----------
\setlist[enumerate]{itemsep=-0.4em,topsep=0.2em}
\setlist[itemize]{itemsep=-0.4em,topsep=0.2em}

% =============================================================================
% Colors — Okabe-Ito colorblind-friendly palette
% =============================================================================

\definecolor{black}{HTML}{000000}
\definecolor{orange}{HTML}{E69F00}
\definecolor{skyblue}{HTML}{56B4E9}
\definecolor{green}{HTML}{009E73}
\definecolor{yellow}{HTML}{F0E442}
\definecolor{blue}{HTML}{0072B2}
\definecolor{vermillion}{HTML}{D55E00}
\definecolor{purple}{HTML}{CC79A7}

% Semantic aliases (mapped to Okabe-Ito)
\definecolor{softbg}{HTML}{F5F5F5}
\definecolor{lightgray}{gray}{0.65}

% =============================================================================
% Box Environments
% =============================================================================

% --- Key points (orange background) ---
\newtcolorbox{keybox}{
  colback=orange!12,
  colframe=orange,
  fonttitle=\bfseries,
  boxrule=0.5pt,
  arc=2pt,
  left=6pt,
  right=6pt,
  top=4pt,
  bottom=4pt
}

% --- Highlights (orange left accent) ---
\newtcolorbox{highlightbox}{
  colback=white,
  colframe=orange,
  fonttitle=\bfseries,
  boxrule=0pt,
  leftrule=3pt,
  arc=0pt,
  left=6pt,
  right=6pt,
  top=4pt,
  bottom=4pt
}

% --- Definitions (blue border, titled) ---
\newtcolorbox{definitionbox}[1][]{
  colback=blue!5,
  colframe=blue,
  fonttitle=\bfseries,
  title={#1},
  boxrule=0.5pt,
  arc=2pt,
  left=6pt,
  right=6pt,
  top=4pt,
  bottom=4pt
}

% --- Methods (sky blue background) ---
\newtcolorbox{methodbox}{
  colback=skyblue!8,
  colframe=skyblue!70,
  fonttitle=\bfseries,
  boxrule=0.5pt,
  arc=2pt,
  left=6pt,
  right=6pt,
  top=4pt,
  bottom=4pt
}

% --- Results (green border) ---
\newtcolorbox{resultbox}{
  colback=green!5,
  colframe=green,
  fonttitle=\bfseries,
  boxrule=0.5pt,
  arc=2pt,
  left=6pt,
  right=6pt,
  top=4pt,
  bottom=4pt
}

% --- Quotes (purple left rule) ---
\newtcolorbox{quotebox}{
  colback=softbg,
  colframe=purple,
  fonttitle=\itshape,
  boxrule=0pt,
  leftrule=2pt,
  arc=0pt,
  left=6pt,
  right=6pt,
  top=4pt,
  bottom=4pt
}

% =============================================================================
% Math Commands — Notation Legend
% =============================================================================

% --- Spaces ---
\newcommand{\R}{\mathbb{R}}
\newcommand{\Rplus}{\mathbb{R}_+}
\newcommand{\Rk}{\mathbb{R}^k}

% --- Probability ---
\newcommand{\Prob}{\mathbb{P}}

% --- Expectation and moments ---
\newcommand{\E}{\mathbb{E}}
\DeclareMathOperator{\var}{var}
\DeclareMathOperator{\cov}{cov}
\DeclareMathOperator{\corr}{corr}
\newcommand{\BLP}{\mathscr{P}}

% --- Convergence ---
\DeclareMathOperator{\plim}{plim}
\newcommand{\convp}{\xrightarrow{\quad p \quad}}
\newcommand{\convd}{\xrightarrow{\quad d \quad}}

% --- Likelihood and information ---
\newcommand{\Lcal}{\mathcal{L}}
\newcommand{\Infmat}{\mathscr{I}}

% --- Relations and operators ---
\newcommand{\defeq}{\overset{\text{def}}{=}}
\newcommand{\ind}{\mathds{1}}
\newcommand{\Normal}{\mathrm{N}}

% --- Calculus shortcuts ---
\newcommand{\suminfty}[1]{\sum_{#1=0}^\infty}
\newcommand{\prtl}[2]{\frac{\partial #1}{\partial #2}}
\renewcommand{\div}[2]{\frac{d #1}{d #2}}
\newcommand{\suchthat}{\text{ s.t. }}
\newcommand{\wrt}{\text{ w.r.t. }}

% --- Text color shortcuts ---
\newcommand{\red}[1]{\textcolor{vermillion}{#1}}
\newcommand{\blue}[1]{\textcolor{blue}{#1}}
\newcommand{\pmark}{\textcolor{green}{\checkmark}}
\newcommand{\xmark}{\textcolor{vermillion}{\texttimes}}

% --- Stata esttab significance stars ---
\newcommand{\sym}[1]{\ensuremath{^{#1}}}

% Compile: cd output/Handouts && TEXINPUTS=../../codes/Preambles:$TEXINPUTS xelatex file.tex
% =============================================================================

% ---------- Packages ----------
% Fonts
\usepackage{palatino}
\usepackage[utf8]{inputenc}
\usepackage[T1]{fontenc}
\usepackage[normalem]{ulem}

% Math
\usepackage{amsmath,amsthm}
\usepackage{newpxmath}
\usepackage{bm}
\usepackage{mathrsfs}
\usepackage{dsfont}

% Figures and tables
\usepackage{graphicx}
\usepackage{placeins}
\usepackage{xcolor}
\usepackage{array,tabularx,booktabs,longtable,subcaption}
\usepackage{threeparttable}
\usepackage{pdflscape}
\usepackage{siunitx}

% TikZ and pgfplots
\usepackage{pgfplots}
\pgfplotsset{compat=1.18}
\usepackage{tikz}
\usetikzlibrary{positioning,arrows.meta,calc}

% Colored boxes
\usepackage[most]{tcolorbox}

% Citations (natbib loaded before hyperref to avoid conflicts)
\usepackage{natbib}

% Miscellaneous
\usepackage{hyperref}
% Links stay clickable but render as plain text: `hidelinks` removes BOTH the
% coloured border boxes (hyperref's default for links/citations/URLs) and any
% link colouring. Must come after the package load.
\hypersetup{hidelinks}
\usepackage{setspace}
\usepackage{enumitem}

% ---------- Margins ----------
\usepackage[left=1in,right=1in,top=0.75in,bottom=1in,marginparwidth=1.5in,footskip=0.7in]{geometry}

% ---------- Itemize ----------
\setlist[enumerate]{itemsep=-0.4em,topsep=0.2em}
\setlist[itemize]{itemsep=-0.4em,topsep=0.2em}

% =============================================================================
% Colors — Okabe-Ito colorblind-friendly palette
% =============================================================================

\definecolor{black}{HTML}{000000}
\definecolor{orange}{HTML}{E69F00}
\definecolor{skyblue}{HTML}{56B4E9}
\definecolor{green}{HTML}{009E73}
\definecolor{yellow}{HTML}{F0E442}
\definecolor{blue}{HTML}{0072B2}
\definecolor{vermillion}{HTML}{D55E00}
\definecolor{purple}{HTML}{CC79A7}

% Semantic aliases (mapped to Okabe-Ito)
\definecolor{softbg}{HTML}{F5F5F5}
\definecolor{lightgray}{gray}{0.65}

% =============================================================================
% Box Environments
% =============================================================================

% --- Key points (orange background) ---
\newtcolorbox{keybox}{
  colback=orange!12,
  colframe=orange,
  fonttitle=\bfseries,
  boxrule=0.5pt,
  arc=2pt,
  left=6pt,
  right=6pt,
  top=4pt,
  bottom=4pt
}

% --- Highlights (orange left accent) ---
\newtcolorbox{highlightbox}{
  colback=white,
  colframe=orange,
  fonttitle=\bfseries,
  boxrule=0pt,
  leftrule=3pt,
  arc=0pt,
  left=6pt,
  right=6pt,
  top=4pt,
  bottom=4pt
}

% --- Definitions (blue border, titled) ---
\newtcolorbox{definitionbox}[1][]{
  colback=blue!5,
  colframe=blue,
  fonttitle=\bfseries,
  title={#1},
  boxrule=0.5pt,
  arc=2pt,
  left=6pt,
  right=6pt,
  top=4pt,
  bottom=4pt
}

% --- Methods (sky blue background) ---
\newtcolorbox{methodbox}{
  colback=skyblue!8,
  colframe=skyblue!70,
  fonttitle=\bfseries,
  boxrule=0.5pt,
  arc=2pt,
  left=6pt,
  right=6pt,
  top=4pt,
  bottom=4pt
}

% --- Results (green border) ---
\newtcolorbox{resultbox}{
  colback=green!5,
  colframe=green,
  fonttitle=\bfseries,
  boxrule=0.5pt,
  arc=2pt,
  left=6pt,
  right=6pt,
  top=4pt,
  bottom=4pt
}

% --- Quotes (purple left rule) ---
\newtcolorbox{quotebox}{
  colback=softbg,
  colframe=purple,
  fonttitle=\itshape,
  boxrule=0pt,
  leftrule=2pt,
  arc=0pt,
  left=6pt,
  right=6pt,
  top=4pt,
  bottom=4pt
}

% =============================================================================
% Math Commands — Notation Legend
% =============================================================================

% --- Spaces ---
\newcommand{\R}{\mathbb{R}}
\newcommand{\Rplus}{\mathbb{R}_+}
\newcommand{\Rk}{\mathbb{R}^k}

% --- Probability ---
\newcommand{\Prob}{\mathbb{P}}

% --- Expectation and moments ---
\newcommand{\E}{\mathbb{E}}
\DeclareMathOperator{\var}{var}
\DeclareMathOperator{\cov}{cov}
\DeclareMathOperator{\corr}{corr}
\newcommand{\BLP}{\mathscr{P}}

% --- Convergence ---
\DeclareMathOperator{\plim}{plim}
\newcommand{\convp}{\xrightarrow{\quad p \quad}}
\newcommand{\convd}{\xrightarrow{\quad d \quad}}

% --- Likelihood and information ---
\newcommand{\Lcal}{\mathcal{L}}
\newcommand{\Infmat}{\mathscr{I}}

% --- Relations and operators ---
\newcommand{\defeq}{\overset{\text{def}}{=}}
\newcommand{\ind}{\mathds{1}}
\newcommand{\Normal}{\mathrm{N}}

% --- Calculus shortcuts ---
\newcommand{\suminfty}[1]{\sum_{#1=0}^\infty}
\newcommand{\prtl}[2]{\frac{\partial #1}{\partial #2}}
\renewcommand{\div}[2]{\frac{d #1}{d #2}}
\newcommand{\suchthat}{\text{ s.t. }}
\newcommand{\wrt}{\text{ w.r.t. }}

% --- Text color shortcuts ---
\newcommand{\red}[1]{\textcolor{vermillion}{#1}}
\newcommand{\blue}[1]{\textcolor{blue}{#1}}
\newcommand{\pmark}{\textcolor{green}{\checkmark}}
\newcommand{\xmark}{\textcolor{vermillion}{\texttimes}}

% --- Stata esttab significance stars ---
\newcommand{\sym}[1]{\ensuremath{^{#1}}}

% Compile: cd output/Handouts && TEXINPUTS=../../codes/Preambles:$TEXINPUTS xelatex file.tex
% =============================================================================

% ---------- Packages ----------
% Fonts
\usepackage{palatino}
\usepackage[utf8]{inputenc}
\usepackage[T1]{fontenc}
\usepackage[normalem]{ulem}

% Math
\usepackage{amsmath,amsthm}
\usepackage{newpxmath}
\usepackage{bm}
\usepackage{mathrsfs}
\usepackage{dsfont}

% Figures and tables
\usepackage{graphicx}
\usepackage{placeins}
\usepackage{xcolor}
\usepackage{array,tabularx,booktabs,longtable,subcaption}
\usepackage{threeparttable}
\usepackage{pdflscape}
\usepackage{siunitx}

% TikZ and pgfplots
\usepackage{pgfplots}
\pgfplotsset{compat=1.18}
\usepackage{tikz}
\usetikzlibrary{positioning,arrows.meta,calc}

% Colored boxes
\usepackage[most]{tcolorbox}

% Citations (natbib loaded before hyperref to avoid conflicts)
\usepackage{natbib}

% Miscellaneous
\usepackage{hyperref}
% Links stay clickable but render as plain text: `hidelinks` removes BOTH the
% coloured border boxes (hyperref's default for links/citations/URLs) and any
% link colouring. Must come after the package load.
\hypersetup{hidelinks}
\usepackage{setspace}
\usepackage{enumitem}

% ---------- Margins ----------
\usepackage[left=1in,right=1in,top=0.75in,bottom=1in,marginparwidth=1.5in,footskip=0.7in]{geometry}

% ---------- Itemize ----------
\setlist[enumerate]{itemsep=-0.4em,topsep=0.2em}
\setlist[itemize]{itemsep=-0.4em,topsep=0.2em}

% =============================================================================
% Colors — Okabe-Ito colorblind-friendly palette
% =============================================================================

\definecolor{black}{HTML}{000000}
\definecolor{orange}{HTML}{E69F00}
\definecolor{skyblue}{HTML}{56B4E9}
\definecolor{green}{HTML}{009E73}
\definecolor{yellow}{HTML}{F0E442}
\definecolor{blue}{HTML}{0072B2}
\definecolor{vermillion}{HTML}{D55E00}
\definecolor{purple}{HTML}{CC79A7}

% Semantic aliases (mapped to Okabe-Ito)
\definecolor{softbg}{HTML}{F5F5F5}
\definecolor{lightgray}{gray}{0.65}

% =============================================================================
% Box Environments
% =============================================================================

% --- Key points (orange background) ---
\newtcolorbox{keybox}{
  colback=orange!12,
  colframe=orange,
  fonttitle=\bfseries,
  boxrule=0.5pt,
  arc=2pt,
  left=6pt,
  right=6pt,
  top=4pt,
  bottom=4pt
}

% --- Highlights (orange left accent) ---
\newtcolorbox{highlightbox}{
  colback=white,
  colframe=orange,
  fonttitle=\bfseries,
  boxrule=0pt,
  leftrule=3pt,
  arc=0pt,
  left=6pt,
  right=6pt,
  top=4pt,
  bottom=4pt
}

% --- Definitions (blue border, titled) ---
\newtcolorbox{definitionbox}[1][]{
  colback=blue!5,
  colframe=blue,
  fonttitle=\bfseries,
  title={#1},
  boxrule=0.5pt,
  arc=2pt,
  left=6pt,
  right=6pt,
  top=4pt,
  bottom=4pt
}

% --- Methods (sky blue background) ---
\newtcolorbox{methodbox}{
  colback=skyblue!8,
  colframe=skyblue!70,
  fonttitle=\bfseries,
  boxrule=0.5pt,
  arc=2pt,
  left=6pt,
  right=6pt,
  top=4pt,
  bottom=4pt
}

% --- Results (green border) ---
\newtcolorbox{resultbox}{
  colback=green!5,
  colframe=green,
  fonttitle=\bfseries,
  boxrule=0.5pt,
  arc=2pt,
  left=6pt,
  right=6pt,
  top=4pt,
  bottom=4pt
}

% --- Quotes (purple left rule) ---
\newtcolorbox{quotebox}{
  colback=softbg,
  colframe=purple,
  fonttitle=\itshape,
  boxrule=0pt,
  leftrule=2pt,
  arc=0pt,
  left=6pt,
  right=6pt,
  top=4pt,
  bottom=4pt
}

% =============================================================================
% Math Commands — Notation Legend
% =============================================================================

% --- Spaces ---
\newcommand{\R}{\mathbb{R}}
\newcommand{\Rplus}{\mathbb{R}_+}
\newcommand{\Rk}{\mathbb{R}^k}

% --- Probability ---
\newcommand{\Prob}{\mathbb{P}}

% --- Expectation and moments ---
\newcommand{\E}{\mathbb{E}}
\DeclareMathOperator{\var}{var}
\DeclareMathOperator{\cov}{cov}
\DeclareMathOperator{\corr}{corr}
\newcommand{\BLP}{\mathscr{P}}

% --- Convergence ---
\DeclareMathOperator{\plim}{plim}
\newcommand{\convp}{\xrightarrow{\quad p \quad}}
\newcommand{\convd}{\xrightarrow{\quad d \quad}}

% --- Likelihood and information ---
\newcommand{\Lcal}{\mathcal{L}}
\newcommand{\Infmat}{\mathscr{I}}

% --- Relations and operators ---
\newcommand{\defeq}{\overset{\text{def}}{=}}
\newcommand{\ind}{\mathds{1}}
\newcommand{\Normal}{\mathrm{N}}

% --- Calculus shortcuts ---
\newcommand{\suminfty}[1]{\sum_{#1=0}^\infty}
\newcommand{\prtl}[2]{\frac{\partial #1}{\partial #2}}
\renewcommand{\div}[2]{\frac{d #1}{d #2}}
\newcommand{\suchthat}{\text{ s.t. }}
\newcommand{\wrt}{\text{ w.r.t. }}

% --- Text color shortcuts ---
\newcommand{\red}[1]{\textcolor{vermillion}{#1}}
\newcommand{\blue}[1]{\textcolor{blue}{#1}}
\newcommand{\pmark}{\textcolor{green}{\checkmark}}
\newcommand{\xmark}{\textcolor{vermillion}{\texttimes}}

% --- Stata esttab significance stars ---
\newcommand{\sym}[1]{\ensuremath{^{#1}}}


\begin{document}

% ─────────────────────────────────────────────────────────────────────────────
\begin{center}
{\Large\bfseries Master's in Asset Management Orientation BootCamp}\\[0.5em]
{\large Probability and Statistics Foundations}\\[0.5em]
% Connection: Pre-Session for MGMT 924 --- Statistical Foundations
{\normalsize Yale School of Management \hspace{1cm} Instructor: Tudor Schlanger}
\end{center}

\vspace{0.5em}
\hrule
\vspace{1em}

% ─────────────────────────────────────────────────────────────────────────────
\section*{Overview}

% Connection: This session provides the probability and statistics foundations needed for MGMT~924 (Statistical Foundations).
This session provides the probability and statistics foundations needed for linear regression. By the end, students will be equipped to understand linear regression as a prediction and inference tool. All concepts are motivated through examples from asset management and financial markets.

\medskip

\noindent\textbf{Prerequisite knowledge:} Basic algebra; no prior statistics required.\\
\textbf{Reference text:} Wooldridge (2016), \textit{Introductory Econometrics}, Appendices B and~C.\\
\textbf{Software:} R (for demonstrations and exercises).

\vspace{1em}

% ═════════════════════════════════════════════════════════════════════════════
\section*{Part 1: Random Variables and Distributions \hfill \normalsize\textnormal{(50 min)}}
% ═════════════════════════════════════════════════════════════════════════════

\noindent\textit{Core question:} How do we describe uncertain outcomes mathematically?

\subsection*{1.1 What is a random variable?}
\begin{itemize}
    \item A random variable as a mapping from outcomes to numbers
    \item Financial examples: tomorrow's stock return, whether a borrower defaults, number of trades in a day
    \item Notation: $X$ for the random variable, $x$ for a realized value
\end{itemize}

\subsection*{1.2 Discrete random variables}
\begin{itemize}
    \item Probability density function (pdf) and how it assigns probabilities to outcomes
    \item Bernoulli distribution --- e.g., does a stock go up or down?
    \item Binomial distribution: number of successes in $n$ independent trials
    \item \textit{Finance example:} Number of days a stock closes higher out of 20 trading days
    \item \textit{Sources:} Wooldridge, Appendix B-1a
\end{itemize}

\subsection*{1.3 Continuous random variables}
\begin{itemize}
    \item Why continuous? Prices, returns, interest rates take on (effectively) infinite values
    \item PDF: probability as area under the curve
    \item Cumulative distribution function (CDF) and its key properties
    \item \textit{Finance example:} What is the probability that a stock loses more than 10\% in a month?
    \item \textit{Sources:} Wooldridge, Appendix B-1b
\end{itemize}

% Connection to MGMT~924: Classes 2--3 assume students can work with random variables, pdf/cdf, and probability statements about stock returns.

\vspace{0.5em}
\hrule height 0.3pt
\vspace{0.5em}
\begin{center}\textit{--- Break (10 min) ---}\end{center}
\vspace{0.5em}
\hrule height 0.3pt
\vspace{1em}

% ═════════════════════════════════════════════════════════════════════════════
\section*{Part 2: Population Parameters and Distribution Moments \hfill \normalsize\textnormal{(50 min)}}
% ═════════════════════════════════════════════════════════════════════════════

\noindent\textit{Core question:} When we know (or assume) the distribution, what summary quantities matter?

\subsection*{2.1 Expected value}
\begin{itemize}
    \item Definition and interpretation: the long-run average; the ``center of mass'' of the distribution
    \item Key properties: linearity of expectation
    \item \textit{Finance example:} Expected return on a portfolio
    \item \textit{Sources:} Wooldridge, Appendix B-3a,b (Properties E.1--E.3)
\end{itemize}

\subsection*{2.2 Variance and standard deviation}
\begin{itemize}
    \item Variance as expected squared deviation from the mean
    \item Standard deviation: same units as $X$
    \item Key properties: effect of linear transformations on variance
    \item Standardization: transforming $X$ to have mean zero and variance one
    \item \textit{Finance example:} Volatility of stock returns; risk measurement
    \item \textit{Sources:} Wooldridge, Appendix B-3d--g (Properties VAR.1--2, SD.1--2)
\end{itemize}

\subsection*{2.3 Covariance and correlation}
\begin{itemize}
    \item Covariance: measuring how two random variables move together
    \item Correlation: standardized covariance, bounded in $[-1, 1]$
    \item Independence implies zero covariance, but not the reverse
    \item Variance of a sum and the role of covariance
    \item \textit{Finance example:} Covariance between stock and bond returns; diversification benefit when $\corr < 1$
    \item \textit{Sources:} Wooldridge, Appendix B-4
\end{itemize}

\subsection*{2.4 The Normal distribution}
\begin{itemize}
    \item Symmetric, bell-shaped, fully characterized by mean and variance
    \item Standard normal and standardization
    \item Linear combinations of normals are normal
    \item 68--95--99.7 rule: probability within 1, 2, 3 standard deviations
    \item \textit{Finance example:} Are stock returns approximately normal? Fat tails and the limits of normality
    \item \textit{Sources:} Wooldridge, Appendix B-5
\end{itemize}

% Connection to MGMT~924: Classes 4--6 (Inference for Linear Regression) rely heavily on the normal distribution, variance decomposition, and the distinction between population and sample moments. Class~12 (Advanced Estimation) revisits moments, skewness, kurtosis, and the method of moments.

\vspace{0.5em}
\hrule height 0.3pt
\vspace{0.5em}
\begin{center}\textit{--- Break (10 min) ---}\end{center}
\vspace{0.5em}
\hrule height 0.3pt
\vspace{1em}

% ═════════════════════════════════════════════════════════════════════════════
\section*{Part 3: Statistical Inference --- Learning from Data \hfill \normalsize\textnormal{(50 min)}}
% ═════════════════════════════════════════════════════════════════════════════

\noindent\textit{Core question:} When we \emph{don't} know the distribution, how do we learn about it from data?

\medskip
\noindent Two fundamental problems:
\begin{enumerate}
    \item \textbf{Estimation:} Learn the value of a population parameter (e.g., $\mu$, $\sigma^2$).
    \item \textbf{Hypothesis testing:} Decide whether we can reject a claim about a population parameter.
\end{enumerate}

\subsection*{3.1 Populations, parameters, and random samples}
\begin{itemize}
    \item Population vs.\ sample; parameter vs.\ statistic
    \item Random sampling: i.i.d.\ draws from a population distribution
    \item \textit{Finance example:} Daily returns on the S\&P 500 as a sample from the population of ``all possible'' trading days
    \item \textit{Sources:} Wooldridge, Appendix C-1
\end{itemize}

\subsection*{3.2 Estimators and estimates}
\begin{itemize}
    \item Statistic = any function of the sample data (e.g., sample mean, sample max)
    \item Estimator = a statistic used to estimate a specific population parameter; estimate = the number you get
    \item The sample mean as an estimator of the population mean
    \item Desirable properties of estimators:
    \begin{itemize}
        \item \textbf{Unbiasedness:} on average, the estimator hits the target
        \item \textbf{Efficiency:} among unbiased estimators, prefer the one with smallest variance
        \item \textbf{Consistency:} as $n$ grows, the estimator converges to the truth
    \end{itemize}
    \item Sample variance and the $n-1$ correction
    \item \textit{Finance example:} Estimating expected return and volatility from historical data
    \item \textit{Sources:} Wooldridge, Appendix C-2 (estimators, unbiasedness, efficiency), C-3 (consistency)
\end{itemize}

\subsection*{3.3 The sampling distribution and the Central Limit Theorem}
\begin{itemize}
    \item The sample mean is itself a random variable with its own distribution (the sampling distribution)
    \item Averaging reduces variability: the variance of the sample mean shrinks with $n$
    \item \textbf{Central Limit Theorem (CLT):} For large $n$, the sample mean is approximately normally distributed regardless of the population distribution
    \item Why the CLT matters: it justifies using normal-based inference even when the population is not normal
    \item \textit{Finance example:} Why portfolio returns look ``more normal'' than individual stock returns (diversification as averaging)
    \item \textit{Sources:} Wooldridge, Appendix C-3
\end{itemize}

\subsection*{3.4 Hypothesis testing}
\begin{itemize}
    \item The framework: null hypothesis $H_0$ vs.\ alternative $H_1$
    \item Test statistics and the $t$-distribution
    \item \textbf{$p$-value:} probability of observing a result at least as extreme as what we got, assuming $H_0$ is true
    \item Decision rule: reject $H_0$ if $p$-value $< \alpha$ (typically $\alpha = 0.05$)
    \item Type~I error (false positive) and Type~II error (false negative)
    \item \textit{Finance example:} Testing whether a trading strategy generates alpha; the multiple testing problem in factor investing
    \item \textit{Sources:} Wooldridge, Appendix C-6
\end{itemize}

\subsection*{3.5 Confidence intervals (if time permits)}
\begin{itemize}
    \item Interval estimation: providing a range of plausible values for a parameter
    \item Interpretation: 95\% of intervals constructed this way contain the true parameter
    \item Relationship to hypothesis testing: $H_0$ is rejected iff the hypothesized value lies outside the CI
    \item \textit{Finance example:} ``We are 95\% confident the expected annual return lies between 4\% and 12\%''
    \item \textit{Sources:} Wooldridge, Appendix C-5
\end{itemize}

% Connection to MGMT~924: These tools are the backbone of Classes 4--6 (Inference for Linear Regression), where students test whether regression coefficients are statistically significant, construct confidence intervals for beta, and use p-values to evaluate hypotheses about asset returns. Class~1 previews the regression setup; this session ensures students arrive with the statistical vocabulary to hit the ground running.

\vspace{1em}

% ─────────────────────────────────────────────────────────────────────────────
% Connection: Mapping to MGMT~924 Curriculum
% This Session                    | MGMT 924 Class | How It Connects
% Random variables, distributions | Class 1        | Understand "stock return" as a random variable
% Expected value, variance        | Classes 2--3   | Prediction = estimating E[Y|X]
% Normal distribution, CLT        | Classes 4--6   | Normality assumption for inference
% Covariance, correlation         | Classes 7--9   | Multiple regression, omitted variable bias
% Estimators, unbiasedness        | Classes 4--6   | OLS as an estimator; Gauss--Markov
% Hypothesis testing, p-values    | Classes 4--6   | Testing H_0: beta = 0 in regression
% Moments (variance, skewness)    | Class 12       | Method of moments estimation
% ─────────────────────────────────────────────────────────────────────────────

\section*{Notation}

See the separate \textit{Notation Legend} document for the complete list of symbols used in this course.

\end{document}
